\documentclass[11pt]{amsart}
\usepackage{amsmath,amssymb,amsthm,mathrsfs}
\usepackage{array,booktabs,tabularx}
\usepackage[margin=1in]{geometry}
\usepackage[colorlinks=true,linkcolor=blue,urlcolor=blue]{hyperref}
\allowdisplaybreaks[2]
\raggedbottom
\newtheorem{theorem}{Theorem}
\newtheorem{proposition}[theorem]{Proposition}
\newtheorem{remark}[theorem]{Remark}
\newcommand{\D}{\mathcal D}
\newcommand{\X}{\mathcal X}
\newcommand{\Y}{\mathcal Y}
\newcommand{\tr}{\operatorname{Tr}}
\newcommand{\supp}{\operatorname{supp}}
\newcommand{\ip}[2]{\langle #1,#2\rangle}
\newcommand{\norm}[1]{\left\|#1\right\|}
\newcommand{\ch}{\operatorname{ch}}
\newcommand{\rank}{\operatorname{rank}}
\newcommand{\fin}{\Subset}
\title{Symbol spaces for bounded and Schatten commutators\\of the downward Haar shift}
\date{2026-11-04}
\begin{document}
\maketitle

\section{Notation and symbol spaces}
We work on $L^2(\mu)$ with the standard dyadic grid $\D$ on
$\mathbb R$. For $I\in\D$, $I_-$ and $I_+$ denote its left and right
children, $\widehat I$ its parent, and $\ell(I)$ its length. We write
$\mathbf1_A$ for the indicator of $A$ and $\#F$ for the cardinality of
a finite set $F$. The notation $\operatorname{Meas}_+(\mathbb R)$ means
the positive Borel measures on $\mathbb R$.
\begin{gather*}
 \D_k=\{[j2^{-k},(j+1)2^{-k}):j\in\mathbb Z\},\qquad
 \D=\bigcup_{k\in\mathbb Z}\D_k,\qquad
 \ch(I)=\{I_-,I_+\},\\
 \mathfrak M=\{\mu\in\operatorname{Meas}_+(\mathbb R):
                   0<\mu(I)<\infty\ (I\in\D)\},\\
 \mathfrak M_\infty=\{\mu\in\mathfrak M:\mu(\mathbb R)=\infty\}
                  \subset\mathfrak M,\qquad \mu\in\mathfrak M.
\end{gather*}
All results below hold for $\mu\in\mathfrak M$ and hence also with
the additional standing assumption $\mu(\mathbb R)=\infty$.
The measure-adapted Haar functions and the downward Haar shift
$H=H_{\D}$ are given by
\begin{gather*}
 m(I)=\frac{\mu(I_-)\mu(I_+)}{\mu(I)},\qquad
 h_I=\sqrt{m(I)}\left(\frac{\mathbf1_{I_-}}{\mu(I_-)}
                     -\frac{\mathbf1_{I_+}}{\mu(I_+)}\right),\\
 \ip{h_I}{h_J}=\delta_{I,J},\qquad
 \sigma_{I_-}=1,\quad\sigma_{I_+}=-1,\qquad
 Hh_I=h_{I_-}-h_{I_+},\quad H^*h_I=\sigma_Ih_{\widehat I}.
\end{gather*}
The inner product is linear in the first variable. The notation
$T^*$ denotes the Hilbert-space adjoint and $T^{\mathsf t}$ the
matrix transpose in the specified orthonormal basis; $\rank T$ is
the dimension of its range. In particular,
\[
 \ip{f}{g}=\int f\overline g\,d\mu,\qquad
 Hf=\sum_{I\in\D}\ip{f}{h_I}(h_{I_-}-h_{I_+}),\qquad f\in L^2(\mu).
\]
The Haar family is completed by constants on dyadic components of
finite measure. We use the following notation for the resulting
orthogonal decomposition and complete basis:
\begin{gather*}
 E_- =(-\infty,0),\qquad E_+=[0,\infty),\qquad
 \mathscr E=\{E\in\{E_-,E_+\}:\mu(E)<\infty\},\\
 \mathcal H_h=\overline{\operatorname{span}}\{h_I:I\in\D\},\qquad
 e_E=\frac{\mathbf1_E}{\sqrt{\mu(E)}}\quad(E\in\mathscr E),\\
 \mathcal H_0=\operatorname{span}\{e_E:E\in\mathscr E\}
             =\mathcal H_h^\perp,\qquad
 L^2(\mu)=\mathcal H_h\oplus\mathcal H_0,\\
 H|_{\mathcal H_0}=H^*|_{\mathcal H_0}=0,\qquad
 H^*H=2P_h,\qquad \norm H=\sqrt2,\qquad P_0=1-P_h,\\
 \mathcal I=\D\sqcup\mathscr E,\qquad
 e_I=h_I\ (I\in\D),\qquad
 \mathscr F=\operatorname{span}\{e_u:u\in\mathcal I\},\\
 F\fin A\ \Longleftrightarrow\ F\subset A,\ \#F<\infty,\qquad
 P_Af=\sum_{u\in A}\ip{f}{e_u}e_u\quad(A\subseteq\mathcal I),\qquad
 \overline{\mathscr F}=L^2(\mu).
\end{gather*}
Here $P_h$ and $P_0$ are the orthogonal projections onto $\mathcal H_h$
and $\mathcal H_0$, respectively. For any set of basis indices $A$,
$P_A$ is the corresponding orthogonal projection. The notation
$F\fin A$ means that $F$ is a finite subset of $A$; $\mathscr F$ is
the space of finite linear combinations of basis vectors. Matrices
act on $\mathscr F$ by the finite sesquilinear pairing
\[
 T=[t_{u,v}]\in\mathbb C^{\mathcal I\times\mathcal I},\quad f,g\in\mathscr F:
 \qquad
 \ip{Tf}{g}:=\sum_{u,v\in\mathcal I}
       t_{u,v}\ip{f}{e_v}\overline{\ip{g}{e_u}}.
\]
Thus $t_{u,v}=\ip{Te_v}{e_u}$ whenever $T$ represents an operator.
We write $\mathcal B(\mathcal H)$ and $\mathcal K(\mathcal H)$ for
the bounded and compact operators on $\mathcal H=L^2(\mu)$.
For $T\in\mathcal K(\mathcal H)$, $s_n(T)$ are the eigenvalues of
$|T|$, listed in nonincreasing order with multiplicity. These are
the singular values of $T$. For $0<p<\infty$,
\[
 |T|=(T^*T)^{1/2},\qquad
 S_p(\mathcal H)=\left\{T\in\mathcal K(\mathcal H):
                      \sum_{n\ge1}s_n(T)^p<\infty\right\},\qquad
 \norm{T}_{S_p}=\left(\sum_{n\ge1}s_n(T)^p\right)^{1/p}.
\]
For $0<p<1$ the expression $\norm{\cdot}_{S_p}$ is a quasi-norm;
for $p\ge1$ it is a norm. The same notation is used for finite matrices.
An unsubscripted operator norm means
$\norm{\cdot}_{L^2(\mu)\to L^2(\mu)}$.
We use $A\lesssim_p B$ and $A\asymp_p B$ for comparison constants
depending only on $p$; without a subscript the constants are absolute.

For a locally integrable complex symbol $\varphi$, $f_I$ denotes a
local average and $\mathbb E_k$ the dyadic conditional expectation.
The coefficients $b_I$ and $d_I$ are, respectively, its Haar
coefficients and the diagonal entries of its multiplication matrix:
\begin{equation}\label{data}
 \begin{gathered}
 \varphi\in L^1_{\mathrm{loc}}(\mu;\mathbb C),\qquad
 f_I=\frac1{\mu(I)}\int_If\,d\mu,\qquad
 \mathbb E_kf=\sum_{I\in\D_k}f_I\mathbf1_I,\\
 b_I=\int\varphi h_I\,d\mu,\\
 d_I=\int\varphi h_I^2\,d\mu
     =\frac{\mu(I_+)\varphi_{I_-}+\mu(I_-)\varphi_{I_+}}{\mu(I)},\\
 d_I-\varphi_I
   =\frac{\mu(I_+)-\mu(I_-)}{\mu(I)}
                     (\varphi_{I_-}-\varphi_{I_+}).
 \end{gathered}
\end{equation}
In particular, $d_I$ is a reverse-weighted average of the two child
means. We write $M_\varphi f=\varphi f$ on its natural domain
$\{f\in L^2(\mu):\varphi f\in L^2(\mu)\}$. For a general locally
integrable symbol the commutator is initially defined by its form on
$\mathscr F$. Membership in an operator class means that this form
has an extension in that class:
\begin{equation}\label{form}
 \begin{gathered}
 c_\varphi(f,g)=\int\varphi\big((Hf)\overline g
                          -f\,\overline{H^*g}\big)\,d\mu,
                          \qquad f,g\in\mathscr F,\\
 C_\varphi\in\mathcal B(L^2(\mu))
 \ \Longleftrightarrow\
 \exists C\in\mathcal B(L^2(\mu))\ 
 \forall f,g\in\mathscr F:\ \ip{Cf}{g}=c_\varphi(f,g),\\
 C_\varphi\in S_p
 \ \Longleftrightarrow\
 \exists C\in S_p\ 
 \forall f,g\in\mathscr F:\ \ip{Cf}{g}=c_\varphi(f,g),\\
 \varphi\in L^\infty(\mu)
 \ \Longrightarrow\
 C_\varphi=[M_\varphi,H]=M_\varphi H-HM_\varphi.
 \end{gathered}
\end{equation}
For each finite $F\subset\D$, $G_F(\varphi)$ records the correlations
of the averaging vectors $u_I$. The trace is unnormalized, and powers
of a positive matrix are defined by spectral calculus.
\begin{equation}\label{gram}
 \begin{gathered}
 u_I=\frac{\mathbf1_I}{\mu(I)},\qquad
 G_F(\varphi)=\left[b_I\overline{b_J}
              \frac{\mu(I\cap J)}{\mu(I)\mu(J)}\right]_{I,J\in F},
              \quad F\fin\D,\\
 z^*G_F(\varphi)z
 =\int\left|\sum_{I\in F}\overline{b_I}z_Iu_I\right|^2d\mu\ge0,\qquad
 \tr A=\sum_{I\in F}A_{I,I}.
 \end{gathered}
\end{equation}
We define the two symbol quantities and their associated homogeneous
space by
\begin{align}
 \mathfrak C_p(\varphi)
   &=\sup_{F\fin\D}\left[\tr\big(G_F(\varphi)^{p/2}\big)\right]^{1/p},
                         \quad 0<p<\infty,\label{carleson}\tag{C}\\
 \mathfrak D_p(\varphi)
   &=\left[\sum_{I\in\D}
       \big(|d_{I_-}-d_I|^2+|d_{I_+}-d_I|^2\big)^{p/2}\right]^{1/p}.
                         \label{gradient}\tag{D}
\end{align}
\begin{gather*}
 \X_p(\mu)=\{\varphi\in L^1_{\mathrm{loc}}(\mu):
       \mathfrak C_p(\varphi)+\mathfrak D_p(\varphi)<\infty\},\qquad
 \norm{\varphi}_{\X_p}=\mathfrak C_p(\varphi)+\mathfrak D_p(\varphi),\\
 \norm{\varphi+c_-\mathbf1_{E_-}+c_+\mathbf1_{E_+}}_{\X_p}
       =\norm{\varphi}_{\X_p},\qquad c_-,c_+\in\mathbb C.
\end{gather*}

The last identity expresses invariance under addition of a constant
on either dyadic component. The quantity $\norm{\cdot}_{\X_p}$ is a
seminorm for $p\ge1$ and a quasi-seminorm for $0<p<1$.

\section{Schatten characterization for \texorpdfstring{$1<p<\infty$}{1 < p < infinity}}
\begin{theorem}\label{main}
Let $\mu\in\mathfrak M$, $1<p<\infty$, and
$\varphi\in L^1_{\mathrm{loc}}(\mu)$. Then
\begin{gather*}
 C_\varphi\in S_p(L^2(\mu))
       \ \Longleftrightarrow\ \varphi\in\X_p(\mu),\\
 c_p\big(\mathfrak C_p(\varphi)+\mathfrak D_p(\varphi)\big)
 \le\norm{C_\varphi}_{S_p}
 \le4\sqrt2\,\mathfrak C_p(\varphi)+\mathfrak D_p(\varphi),\qquad
 c_p>0.
\end{gather*}
Here $c_p$ depends only on $p$.
\end{theorem}
\begin{proof}
Define the finite-row paraproducts $\pi_{\varphi,F}$ and the
paraproduct matrix $\pi_\varphi$ by
\begin{gather*}
 \pi_{\varphi,F}f=\sum_{I\in F}b_I\ip{f}{u_I}h_I,\qquad
 [\pi_{\varphi,F}\pi_{\varphi,F}^*]_{F\times F}=G_F(\varphi),\\
 \norm{\pi_{\varphi,F}}_{S_p}^p
        =\tr\big(G_F(\varphi)^{p/2}\big),\qquad
 (\pi_\varphi f)_I=b_I\ip{f}{u_I}.
\end{gather*}
\begin{equation}\label{pi-norm}
 \sup_{F\fin\D}\norm{\pi_{\varphi,F}}_{S_p}<\infty
 \ \Longleftrightarrow\ \pi_\varphi\in S_p,\qquad
 \norm{\pi_\varphi}_{S_p}=\mathfrak C_p(\varphi).
\end{equation}
Let $D_\varphi$ be the diagonal part and $\rho_I$ its edge difference:
\begin{gather*}
 D_\varphi h_I=d_Ih_I,\qquad D_\varphi|_{\mathcal H_0}=0,\qquad
 \rho_I=\big(|d_{I_-}-d_I|^2+|d_{I_+}-d_I|^2\big)^{1/2},\\
 [D_\varphi,H]h_I
       =(d_{I_-}-d_I)h_{I_-}-(d_{I_+}-d_I)h_{I_+},\\
 \ip{[D_\varphi,H]h_I}{[D_\varphi,H]h_J}
                  =\delta_{I,J}\rho_I^2.
\end{gather*}
\begin{equation}\label{D-norm}
 \norm{[D_\varphi,H]}_{S_p}^p=\sum_I\rho_I^p
                         =\mathfrak D_p(\varphi)^p.
\end{equation}
For $I\subsetneq J$, $h_J(I)$ denotes the constant value of $h_J$ on
$I$. The multiplication matrix and the resulting form decomposition
are
\begin{equation}\label{multiplication-matrix}
 \begin{gathered}
 I\subsetneq J\ \Longrightarrow\ h_J|_I\equiv h_J(I),\\
 a_{I,J}:=\int\varphi h_Ih_J\,d\mu=
 \begin{cases}
 d_I,&I=J,\\
 b_I h_J(I),&I\subsetneq J,\\
 b_J h_I(J),&J\subsetneq I,\\
 0,&I\cap J=\varnothing.
 \end{cases}
 \end{gathered}
\end{equation}
All matrix identities here are understood on $\mathscr F$ before the
corresponding operator extensions are established.
\begin{equation}\label{decompose}
 \begin{gathered}
 c_\varphi(f,g)
  =\ip{[\pi_\varphi,H]f}{g}
    +\ip{[(\pi_{\overline\varphi})^*,H]f}{g}
    +\ip{[D_\varphi,H]f}{g},\quad f,g\in\mathscr F,\\
 \pi_\varphi\in S_p,\ [D_\varphi,H]\in S_p
 \ \Longrightarrow\
 C_\varphi=[\pi_\varphi,H]+[(\pi_{\overline\varphi})^*,H]
                                    +[D_\varphi,H]\in S_p.
 \end{gathered}
\end{equation}
\emph{Sufficiency.}
\begin{align*}
 \varphi\in\X_p
 &\Longrightarrow
 \norm{\pi_{\overline\varphi}}_{S_p}
        =\norm{\pi_\varphi}_{S_p}=\mathfrak C_p(\varphi),\\
 \norm{C_\varphi}_{S_p}
 &\le2\norm H\big(\norm{\pi_\varphi}_{S_p}
                      +\norm{\pi_{\overline\varphi}}_{S_p}\big)
                       +\norm{[D_\varphi,H]}_{S_p}\\
 &=4\sqrt2\,\mathfrak C_p(\varphi)+\mathfrak D_p(\varphi).
\end{align*}
\emph{Necessity.} Denote by $C_{hh}$ the Haar compression of $C$,
and by $\mathscr R$ the Schur multiplier retaining disjoint intervals.
The triangular-projection bound gives
\begin{gather*}
 C=C_\varphi\in S_p,\qquad C_{hh}=P_hC|_{\mathcal H_h},\qquad
 (\mathscr RT)_{I,J}=\mathbf1_{\{I\cap J=\varnothing\}}T_{I,J},\\
 I=[\ell_I,r_I),\quad J=[\ell_J,r_J)
 \ \Longrightarrow\
 \mathbf1_{\{I\cap J=\varnothing\}}
   =\mathbf1_{\{r_I\le\ell_J\}}
      +\mathbf1_{\{r_J\le\ell_I\}},\\
 1<p<\infty
 \ \overset{\text{\cite[\S 2.4]{PST}}}{\Longrightarrow}\
 \norm{\mathscr RT}_{S_p}\le\kappa_p\norm T_{S_p},\qquad
 0<\kappa_p<\infty.
\end{gather*}
Let $U$ be the strict ancestor part of the multiplication matrix.
\begin{gather*}
 U_{I,J}=\mathbf1_{\{J\subsetneq I\}}b_Jh_I(J),\qquad
 C_{K,J}=a_{K,J_-}-a_{K,J_+}-\sigma_K a_{\widehat K,J},\\
 K\cap J=\varnothing
 \ \Longrightarrow\ C_{K,J}=-\sigma_Ka_{\widehat K,J},\\
 J\subsetneq I
 \ \Longrightarrow\
 \#\{K\in\ch(I):K\cap J=\varnothing\}=1,\\
 \sum_{K\in\ch(I)}\sigma_K
       \mathbf1_{\{K\cap J=\varnothing\}}C_{K,J}
      =-\mathbf1_{\{J\subsetneq I\}}b_Jh_I(J).
\end{gather*}
\begin{equation}\label{recovery}
 U=-H^*\mathscr R(C_{hh}),\qquad
 P_h\pi_\varphi|_{\mathcal H_h}=U^{\mathsf t}.
\end{equation}
The remaining paraproduct columns are indexed by $\mathcal H_0$:
\begin{gather*}
 \norm{U^{\mathsf t}}_{S_p}=\norm U_{S_p}
       \le\sqrt2\kappa_p\norm C_{S_p},\\
 B=\pi_\varphi P_0,\qquad
 P_hCP_0=-HB,\qquad B=-\tfrac12H^*P_hCP_0,\\
 \norm B_{S_p}\le2^{-1/2}\norm C_{S_p},\qquad
 \mathfrak C_p(\varphi)=\norm{\pi_\varphi}_{S_p}
       \le(\sqrt2\kappa_p+2^{-1/2})\norm C_{S_p}.
\end{gather*}
Finally, $P_k$ projects onto the Haar functions at level $k$, and
$\mathscr E_1$ extracts one-generation matrix entries:
\begin{gather*}
 P_kf=\sum_{I\in\D_k}\ip{f}{h_I}h_I,\qquad
 V_th_I=e^{ikt}h_I\quad(I\in\D_k),\\
 \mathscr E_1(T)=\sum_kP_{k+1}TP_k
   =\frac1{2\pi}\int_0^{2\pi}e^{-it}V_tTV_t^*\,dt,\\
 \norm{\mathscr E_1(T)}_{S_p}
 \le\frac1{2\pi}\int_0^{2\pi}\norm{V_tTV_t^*}_{S_p}\,dt
 =\norm T_{S_p},\\
 C_{I_-,I}=d_{I_-}-d_I,\qquad
 C_{I_+,I}=-(d_{I_+}-d_I),\qquad
 \mathscr E_1(C_{hh})=[D_\varphi,H],\\
 \mathfrak D_p(\varphi)\le\norm C_{S_p},\qquad
 \norm{\varphi}_{\X_p}
 \le(\sqrt2\kappa_p+2^{-1/2}+1)\norm C_{S_p}.
\end{gather*}
\end{proof}

\begin{proposition}[$p=2$]\label{p2}
Under the standing measure assumptions,
\[
 \norm{C_\varphi}_{S_2}^2\asymp
 \sum_{I\in\D}\frac{|b_I|^2}{\mu(I)}
 +\sum_{I\in\D}\big(|d_{I_-}-d_I|^2+|d_{I_+}-d_I|^2\big).
\]
\end{proposition}
\begin{proof}
\begin{gather*}
 \tr G_F=\sum_{I\in F}\frac{|b_I|^2}{\mu(I)}
 \ \Longrightarrow\
 \mathfrak C_2(\varphi)^2
       =\sum_{I\in\D}\frac{|b_I|^2}{\mu(I)},\\
 \eqref{D-norm},\quad(\ref{main})
 \ \Longrightarrow\
 \norm{C_\varphi}_{S_2}^2
       \asymp\mathfrak C_2(\varphi)^2+\mathfrak D_2(\varphi)^2.
\end{gather*}
\end{proof}

\section{Finite-matrix characterization for \texorpdfstring{$0<p<\infty$}{0 < p < infinity}}
Let $\gamma_{u,v}$ denote the entries of the commutator form in the
complete basis $\{e_u\}_{u\in\mathcal I}$. They are given by
\begin{equation}\label{gamma}
 \gamma_{I,J}=a_{I,J_-}-a_{I,J_+}-\sigma_Ia_{\widehat I,J},
                      \qquad I,J\in\D.
\end{equation}
\begin{gather*}
 E,E'\in\mathscr E,\quad I,J\in\D
 \ \Longrightarrow\\
 \gamma_{I,E}=-\frac{\sigma_Ib_{\widehat I}}{\sqrt{\mu(E)}}
                       \mathbf1_{\{I\subset E\}},\qquad
 \gamma_{E,J}=\frac{b_{J_-}-b_{J_+}}{\sqrt{\mu(E)}}
                       \mathbf1_{\{J\subset E\}},\qquad
 \gamma_{E,E'}=0,\\
 \gamma_{u,v}=c_\varphi(e_v,e_u),\qquad
 \Gamma_F=[\gamma_{u,v}]_{u,v\in F},\qquad F\fin\mathcal I.
\end{gather*}
Define the homogeneous space $\Y_p(\mu)$ from these finite sections:
\begin{equation}\label{Y-norm}
 \begin{gathered}
 \norm{\varphi}_{\Y_p}
  =\sup_{F\fin\mathcal I}
       \left[\tr\big((\Gamma_F^*\Gamma_F)^{p/2}\big)\right]^{1/p},\\
 \Y_p(\mu)=\{\varphi\in L^1_{\mathrm{loc}}(\mu):
                                      \norm{\varphi}_{\Y_p}<\infty\}.
 \end{gathered}
\end{equation}
This construction retains the full commutator matrix. For
$1<p<\infty$, Theorem~\ref{main} separates its condition into
$\mathfrak C_p$ and $\mathfrak D_p$.
\begin{theorem}\label{all-p}
For $\mu\in\mathfrak M$ and $\varphi\in L^1_{\mathrm{loc}}(\mu)$,
\begin{gather*}
 0<p<\infty:\qquad
 C_\varphi\in S_p(L^2(\mu))\ \Longleftrightarrow\ \varphi\in\Y_p(\mu),
 \qquad\norm{C_\varphi}_{S_p}=\norm{\varphi}_{\Y_p},\\
 1<p<\infty:\qquad
 \Y_p(\mu)=\X_p(\mu),\qquad
 \norm{\varphi}_{\Y_p}\asymp_p\norm{\varphi}_{\X_p}.
\end{gather*}
\end{theorem}
\begin{proof}
The finite-section bound first determines a bounded extension;
the Schatten Fatou property then gives its class membership:
\begin{gather*}
 M:=\sup_{F\fin\mathcal I}\norm{\Gamma_F}_{S_p}<\infty,\qquad
 \norm{\Gamma_F}_{\ell^2(F)\to\ell^2(F)}\le\norm{\Gamma_F}_{S_p},\\
 f,g\in\mathscr F,\quad f=P_Ff,\quad g=P_Fg
 \ \Longrightarrow\
 |c_\varphi(f,g)|\le M\norm f_2\norm g_2,\\
 \overline{\mathscr F}=L^2(\mu)
 \ \Longrightarrow\
 \exists!\, C\in\mathcal B(L^2(\mu)):\
 \ip{Cf}{g}=c_\varphi(f,g),\quad \norm C\le M,\\
 P_FCP_F=\Gamma_F,\qquad
 \sup_{F\fin\mathcal I}\norm{P_FCP_F}_{S_p}=M
 \ \overset{\mathrm{Fatou}}{\Longrightarrow}\
 C\in S_p,\quad\norm C_{S_p}\le M.
\end{gather*}
\begin{gather*}
 C=C_\varphi\in S_p,\quad F_n\fin\mathcal I,\quad F_n\uparrow\mathcal I
 \ \Longrightarrow\
 \norm{P_{F_n}CP_{F_n}-C}_{S_p}\longrightarrow0,\\
 \norm C_{S_p}
  =\lim_n\norm{\Gamma_{F_n}}_{S_p}
  \le\sup_{F\fin\mathcal I}\norm{\Gamma_F}_{S_p}
  \le\norm C_{S_p},\\
 1<p<\infty,\quad(\ref{main})
 \ \Longrightarrow\
 \Y_p=\X_p,\quad\norm{\varphi}_{\Y_p}\asymp_p\norm{\varphi}_{\X_p}.
\end{gather*}
\end{proof}

\section{An endpoint counterexample for \texorpdfstring{$0<p\le1$}{0 < p <= 1}}\label{sec-endpoint-example}
The following measure and rank-one multiplier yield a symbol in
$\Y_p(\mu)\setminus\X_p(\mu)$ for every $0<p\le1$. Here $dx$ is
Lebesgue measure, $\delta_0$ is the unit point mass at $0$, and $|J|$
denotes the Lebesgue length of $J$.
\begin{gather*}
 I_n=[0,2^{-n}),\quad J_n=I_n\setminus I_{n+1},\quad n\in\mathbb N_0,\\
 N_k=2^k,\qquad n_0=0,\quad n_k=\sum_{r=1}^k2^r=2^{k+1}-2,\qquad
 B_k=\{n_{k-1},\ldots,n_k-1\},\\
 \mathbb N_0=\bigsqcup_{k\ge1}B_k,\quad\#B_k=N_k,\quad
 w_k=\frac3{\pi^2k^2},\quad
 a_n=\frac{w_k}{N_k}\quad(n\in B_k),\\
 s_0=\tfrac12,\qquad s_{n+1}=s_n+a_n,\qquad
 \sum_na_n=\sum_kw_k=\tfrac12,\qquad s_n\uparrow1,\\
 d\mu=\delta_0+\mathbf1_{\mathbb R\setminus I_0}\,dx
       +\sum_{n\ge0}\frac{s_n^{-1}-s_{n+1}^{-1}}{|J_n|}
                                      \mathbf1_{J_n}\,dx,\\
 \mu\in\mathfrak M_\infty,\qquad
 \mu(\{0\})=1,\quad\mu(I_n)=s_n^{-1},\qquad
 \varphi=\mathbf1_{\{0\}},\\
 M_\varphi f=f(0)\mathbf1_{\{0\}},\qquad
 \rank M_\varphi=1
 \ \Longrightarrow\
 \rank C_\varphi\le2
 \ \Longrightarrow\
 \varphi\in\bigcap_{p>0}\Y_p(\mu).
\end{gather*}
The symbol $Q_k$ below denotes the orthogonal projection onto the
Haar indices in the block $B_k$; $L_N$ is a strict lower-triangular
matrix.
\begin{gather*}
 b_{I_n}=h_{I_n}(0)
 =\sqrt{\mu(I_{n+1})^{-1}-\mu(I_n)^{-1}}=\sqrt{a_n},\\
 \ip{\pi_\varphi h_{I_j}}{h_{I_i}}
   =\begin{cases}\sqrt{a_ia_j},&i>j,\\0,&i\le j,\end{cases}\qquad
 Q_kf=\sum_{n\in B_k}\ip{f}{h_{I_n}}h_{I_n},\\
 Q_k\pi_\varphi Q_k=\frac{w_k}{N_k}L_{N_k},\qquad
 (L_N)_{i,j}=\mathbf1_{\{i>j\}},\\
 s_j(L_N)=\left[2\sin\left(\frac{(2j-1)\pi}{4N-2}\right)\right]^{-1},
                         \quad1\le j\le N-1,\\
 \norm{L_N}_{S_1}
  \ge\frac{2N-1}{\pi}\sum_{j=1}^{N-1}\frac1{2j-1}
  \gtrsim N\log N.
\end{gather*}
Block pinching in trace norm gives
\begin{align*}
 \pi_\varphi\in S_1
 &\Longrightarrow
 \norm{\pi_\varphi}_{S_1}
   \ge\sum_{k=1}^m\norm{Q_k\pi_\varphi Q_k}_{S_1}\\
 &\gtrsim\sum_{k=1}^m w_k\log N_k
   \asymp\sum_{k=1}^m\frac1k\xrightarrow[m\to\infty]{}\infty
 \ \Longrightarrow\ \bot.
\end{align*}
\begin{gather*}
 \pi_\varphi\notin S_1,\qquad
 0<p\le1\ \Longrightarrow\ S_p\subset S_1
 \ \Longrightarrow\
 \mathfrak C_p(\varphi)=\infty,\\
 \varphi\in\bigcap_{0<p\le1}\big(\Y_p(\mu)\setminus\X_p(\mu)\big).
\end{gather*}

\section{Boundedness on \texorpdfstring{$L^2(\mu)$}{L2(mu)}}
The operator-norm endpoint is denoted by $p=\infty$. Define
$\mathfrak A(\varphi)$ as the local $L^2$ mean-oscillation seminorm
and $\mathfrak D_\infty(\varphi)$ as the uniform edge seminorm of the
diagonal coefficients $d_I$:
\begin{align}
 \mathfrak A(\varphi)
  &=\sup_{Q\in\D}\left(\frac1{\mu(Q)}
                \int_Q|\varphi-\varphi_Q|^2\,d\mu\right)^{1/2}
    =\norm{\varphi}_{\mathrm{bmo}_{\mu,2}},\label{local-bmo}\\
 \mathfrak D_\infty(\varphi)
  &=\sup_{I\in\D}
       \big(|d_{I_-}-d_I|^2+|d_{I_+}-d_I|^2\big)^{1/2}.
       \label{edge-infty}
\end{align}
\begin{gather*}
 \X_\infty(\mu)=\{\varphi\in L^1_{\mathrm{loc}}(\mu):
       \mathfrak A(\varphi)+\mathfrak D_\infty(\varphi)<\infty\},\\
 \norm{\varphi}_{\X_\infty}
          =\mathfrak A(\varphi)+\mathfrak D_\infty(\varphi).
\end{gather*}
\begin{theorem}\label{boundedness}
Let $\mu\in\mathfrak M$ and $\varphi\in L^1_{\mathrm{loc}}(\mu)$. Then
\begin{gather*}
 C_\varphi\in\mathcal B(L^2(\mu))
       \ \Longleftrightarrow\ \varphi\in\X_\infty(\mu).
\end{gather*}
Whenever these conditions hold, write $C=C_\varphi$. One has
\begin{align}
 \mathfrak A(\varphi)&\le\sqrt{3/2}\,\norm C,
                                          \label{A-lower}\\
 \mathfrak D_\infty(\varphi)&\le\norm C,
                                          \label{D-lower}\\
 \norm C&\le8\sqrt2\,\mathfrak A(\varphi)+\mathfrak D_\infty(\varphi).
                                          \label{bounded-upper}
\end{align}
\begin{gather*}
 \mathscr E=\varnothing\ \Longrightarrow\ \mathfrak A(\varphi)\le\norm C,\\
 \norm C\asymp\norm{\varphi}_{\mathrm{bmo}_{\mu,2}}
                      +\sup_{I\in\D}|d_I-d_{\widehat I}|.
\end{gather*}
\end{theorem}
\begin{proof}
First, local martingale approximation gives the coefficient identity.
For $Q\in\D$, $k_Q$ is determined by $Q\in\D_{k_Q}$.
\begin{gather*}
 Q\in\D_{k_Q},\qquad
 S_{Q,n}=\sum_{\substack{J\subseteq Q\\
             0\le\log_2(\ell(Q)/\ell(J))<n}}b_Jh_J
       =(\mathbb E_{k_Q+n}\varphi-\varphi_Q)\mathbf1_Q,\\
 S_{Q,n}\xrightarrow[n\to\infty]{L^1(Q,\mu)}
                       (\varphi-\varphi_Q)\mathbf1_Q,\qquad
 \norm{S_{Q,n}}_2^2
       =\sum_{\substack{J\subseteq Q\\
             0\le\log_2(\ell(Q)/\ell(J))<n}}|b_J|^2,\\
 \sum_{J\subseteq Q}|b_J|^2<\infty
 \ \Longleftrightarrow\
 (\varphi-\varphi_Q)\mathbf1_Q\in L^2(\mu),\\
 \int_Q|\varphi-\varphi_Q|^2d\mu
                   =\sum_{J\subseteq Q}|b_J|^2\in[0,\infty].
\end{gather*}
\begin{equation}\label{carleson-local}
 \mathfrak A(\varphi)^2
   =\sup_{Q\in\D}\frac1{\mu(Q)}\sum_{J\subseteq Q}|b_J|^2.
\end{equation}
\emph{Necessity.} Assume $C$ is bounded. Its child entries control
the diagonal edge seminorm:
\begin{gather*}
 C=C_\varphi\in\mathcal B(L^2(\mu)),\qquad N=\norm C,\\
 C_{I_-,I}=d_{I_-}-d_I,\qquad
 C_{I_+,I}=-(d_{I_+}-d_I),\\
 \rho_I^2=|C_{I_-,I}|^2+|C_{I_+,I}|^2
       \le\norm{Ch_I}_2^2\le N^2
 \ \Longrightarrow\ \mathfrak D_\infty(\varphi)\le N.
\end{gather*}
Fix a dyadic interval $Q$ in the component $E$. For each strict
ancestor $A$ of $Q$, let $A_Q$ be the child containing $Q$ and $K_A$
the other child.
\begin{gather*}
 Q\subset E\in\{E_-,E_+\},\qquad A\in\D,\quad A\supsetneq Q,\\
 A_Q\in\ch(A),\quad Q\subseteq A_Q,\qquad
 K_A\in\ch(A)\setminus\{A_Q\},\qquad K_A\cap Q=\varnothing,\\
 J\subseteq Q\ \Longrightarrow\
 \supp Hh_J\subseteq J,\quad
 \ip{\varphi Hh_J}{h_{K_A}}=0,\quad
 \int\varphi h_JH^*h_{K_A}\,d\mu
                  =\sigma_{K_A}h_A(Q)b_J.
\end{gather*}
\begin{equation}\label{outside-entry}
 C_{K_A,J}=-\sigma_{K_A}h_A(Q)b_J,
                 \qquad A\supsetneq Q,\quad J\subseteq Q.
\end{equation}
\begin{gather*}
 |h_A(Q)|^2=\frac1{\mu(A_Q)}-\frac1{\mu(A)},\qquad
 A\ne B\ \Longrightarrow\ K_A\cap K_B=\varnothing,\\
 Q^{(0)}=Q,\qquad Q^{(n+1)}=\widehat{Q^{(n)}},\qquad
 Q^{(n)}\uparrow E,\qquad \mu(Q^{(n)})\uparrow\mu(E).
\end{gather*}
\begin{equation}\label{telescoping-exterior}
 \sum_{A\supsetneq Q}|h_A(Q)|^2
  =\lim_{n\to\infty}\sum_{j=1}^n
       \left(\frac1{\mu(Q^{(j-1)})}-\frac1{\mu(Q^{(j)})}\right)
  =\frac1{\mu(Q)}-\frac1{\mu(E)},\qquad \frac1\infty=0.
\end{equation}
Write $\D(Q)$ for the descendants of $Q$, including $Q$, and
$\D^{\mathrm{out}}(Q)$ for the dyadic intervals disjoint from $Q$.
The exterior compression has the following rank-one form:
\begin{gather*}
 \D(Q)=\{J\in\D:J\subseteq Q\},\qquad
 \D^{\mathrm{out}}(Q)=\{K\in\D:K\cap Q=\varnothing\},\\
 v_Q=\sum_{A\supsetneq Q}\sigma_{K_A}h_A(Q)h_{K_A},\qquad
 F\fin\D(Q),\\
 P_{\D^{\mathrm{out}}(Q)}CP_Ff
          =-\left(\sum_{J\in F}b_J\ip{f}{h_J}\right)v_Q,\\
 \norm{P_{\D^{\mathrm{out}}(Q)}CP_F}^{2}
   =\norm{v_Q}_2^2\sum_{J\in F}|b_J|^2
   =\left(\frac1{\mu(Q)}-\frac1{\mu(E)}\right)
                          \sum_{J\in F}|b_J|^2\le N^2.
\end{gather*}
\begin{equation}\label{exterior-carleson}
 F\uparrow\D(Q)\ \Longrightarrow\
 \left(\frac1{\mu(Q)}-\frac1{\mu(E)}\right)
                   \sum_{J\subseteq Q}|b_J|^2\le N^2.
\end{equation}
\begin{gather*}
 \mu(E)=\infty\ \Longrightarrow\
 \frac1{\mu(Q)}\sum_{J\subseteq Q}|b_J|^2\le N^2,\\
 E\in\mathscr E\ \Longrightarrow\
 \ip{Ce_E}{h_I}
    =-\frac{\sigma_I b_{\widehat I}}{\sqrt{\mu(E)}}
                                    \mathbf1_{\{I\subset E\}},\quad
 P_0Ce_E=0.
\end{gather*}
For a component of finite measure, its constant basis vector supplies
the remaining Carleson contribution:
\begin{equation}\label{finite-component-column}
 E\in\mathscr E\ \Longrightarrow\
 \norm{Ce_E}_2^2
  =\frac1{\mu(E)}\sum_{I\subset E}|b_{\widehat I}|^2
  =\frac2{\mu(E)}\sum_{J\subset E}|b_J|^2\le N^2.
\end{equation}
\begin{gather*}
 E\in\mathscr E,\quad Q\subset E
 \ \Longrightarrow\\
 \frac1{\mu(Q)}\sum_{J\subseteq Q}|b_J|^2
  \le N^2+\frac1{\mu(E)}\sum_{J\subset E}|b_J|^2
  \le\frac32N^2,\\
 \eqref{carleson-local}
 \ \Longrightarrow\
 \mathfrak A(\varphi)\le\sqrt{3/2}\,N,\qquad
 \mathscr E=\varnothing\ \Longrightarrow\ \mathfrak A(\varphi)\le N.
\end{gather*}
\emph{Sufficiency.} Let $M_{\D}$ be the dyadic maximal operator.
For finite $F$ and $t>0$, $\mathcal M_t(F)$ consists of the
inclusion-maximal intervals in $\{I\in F:|f_I|>t\}$. Doob's $L^2$
inequality and maximal selection give
\begin{gather*}
 \varphi\in\X_\infty,\quad f\in L^2(\mu),\quad
 M_{\D}f(x)=\sup_{I\ni x}\frac1{\mu(I)}\int_I|f|\,d\mu,\qquad
 \norm{M_{\D}f}_2\le2\norm f_2,\\
 F\fin\D,\quad t>0,\qquad
 \mathcal M_t(F)=\max_{\subseteq}\{I\in F:|f_I|>t\},\\
 Q,R\in\mathcal M_t(F),\quad Q\ne R
 \ \Longrightarrow\ Q\cap R=\varnothing,\qquad
 \bigcup_{Q\in\mathcal M_t(F)}Q\subseteq\{M_{\D}f>t\},\\
 \sum_{\substack{I\in F\\|f_I|>t}}|b_I|^2
  \le\sum_{Q\in\mathcal M_t(F)}\sum_{I\subseteq Q}|b_I|^2
  \le\mathfrak A(\varphi)^2\sum_{Q\in\mathcal M_t(F)}\mu(Q)
  \le\mathfrak A(\varphi)^2\mu\{M_{\D}f>t\}.
\end{gather*}
\begin{align*}
 \norm{\pi_{\varphi,F}f}_2^2
  &=\sum_{I\in F}|b_I|^2|f_I|^2
   =\int_0^\infty2t\sum_{\substack{I\in F\\|f_I|>t}}|b_I|^2\,dt\\
  &\le\mathfrak A(\varphi)^2
           \int_0^\infty2t\,\mu\{M_{\D}f>t\}\,dt\\
  &=\mathfrak A(\varphi)^2\norm{M_{\D}f}_2^2
   \le4\mathfrak A(\varphi)^2\norm f_2^2.
\end{align*}
\begin{equation}\label{pi-bounded}
 \sup_{F\fin\D}\norm{\pi_{\varphi,F}}\le2\mathfrak A(\varphi)
 \ \Longrightarrow\
 \norm{\pi_\varphi}=\norm{\pi_{\overline\varphi}}
                         \le2\mathfrak A(\varphi).
\end{equation}
The diagonal commutator and the form decomposition now give
\begin{gather*}
 \norm{[D_\varphi,H]f}_2^2
       =\sum_{I\in\D}\rho_I^2|\ip{f}{h_I}|^2
 \ \Longrightarrow\
 \norm{[D_\varphi,H]}=\sup_I\rho_I=\mathfrak D_\infty(\varphi),\\
 \eqref{decompose},\quad\eqref{pi-bounded}
 \ \Longrightarrow\
 \norm{C_\varphi}
  \le2\norm H\big(\norm{\pi_\varphi}+\norm{\pi_{\overline\varphi}}\big)
                       +\mathfrak D_\infty(\varphi)
  \le8\sqrt2\,\mathfrak A(\varphi)+\mathfrak D_\infty(\varphi),\\
 \sup_I|d_I-d_{\widehat I}|\le\mathfrak D_\infty(\varphi)
                    \le\sqrt2\sup_I|d_I-d_{\widehat I}|.
\end{gather*}
\end{proof}

\begin{remark}[Paraproduct norm]
Testing on normalized interval indicators gives
\begin{gather*}
 g_Q=\frac{\mathbf1_Q}{\sqrt{\mu(Q)}},\qquad\norm{g_Q}_2=1,\\
 \norm{\pi_\varphi g_Q}_2^2
       \ge\frac1{\mu(Q)}\sum_{J\subseteq Q}|b_J|^2,\\
 \mathfrak A(\varphi)\le\norm{\pi_\varphi}
                         \le2\mathfrak A(\varphi)
 \ \Longrightarrow\\
 \norm{C_\varphi}\asymp\norm{\pi_\varphi}
                            +\mathfrak D_\infty(\varphi).
\end{gather*}
\end{remark}

\begin{remark}[Lebesgue measure]
For Lebesgue measure, the diagonal edge condition is controlled by
the dyadic mean-oscillation seminorm:
\begin{gather*}
 d\mu=dx\ \Longrightarrow\
 d_I=\varphi_I,\qquad
 d_{I_-}-d_I=\frac{b_I}{\sqrt{\mu(I)}},\quad
 d_{I_+}-d_I=-\frac{b_I}{\sqrt{\mu(I)}},\\
 \mathfrak D_\infty(\varphi)
     =\sqrt2\sup_I\frac{|b_I|}{\sqrt{\mu(I)}}
     \le\sqrt2\mathfrak A(\varphi),\\
 \norm{C_\varphi}\asymp\mathfrak A(\varphi)
                  =\norm{\varphi}_{\mathrm{BMO}_2^{\D}}.
\end{gather*}
\end{remark}

\begin{remark}[The additional diagonal condition]
The following atomless example has finite local mean oscillation
and an unbounded diagonal edge seminorm. Each unit interval has
measure one, while its two halves may have very different masses.
\begin{gather*}
 Q_n=[n,n+1),\quad\varepsilon_n=2^{-2n},\quad n\ge1,\\
 d\mu=\mathbf1_{\mathbb R\setminus[1,\infty)}\,dx
      +\sum_{n\ge1}\big(2\varepsilon_n\mathbf1_{(Q_n)_-}
              +2(1-\varepsilon_n)\mathbf1_{(Q_n)_+}\big)\,dx,\\
 \mu\in\mathfrak M_\infty,\qquad\mu(\{x\})=0\quad(x\in\mathbb R),
 \qquad \mu([j,j+1))=1\quad(j\in\mathbb Z),\\
 \varphi=\sum_{n\ge1}\varepsilon_n^{-1/2}\mathbf1_{(Q_n)_-}.
\end{gather*}
\begin{gather*}
 I\in\D,\quad\ell(I)<1
 \ \Longrightarrow\
 \varphi|_I\equiv\varphi_I
 \ \Longrightarrow\ \int_I|\varphi-\varphi_I|^2d\mu=0,\\
 I\in\D,\quad\ell(I)\ge1
 \ \Longrightarrow\\
 \frac1{\mu(I)}\int_I|\varphi-\varphi_I|^2d\mu
 \le\frac1{\mu(I)}\int_I|\varphi|^2d\mu
 =\frac{\#\{n\ge1:Q_n\subseteq I\}}{\mu(I)}\le1,\\
 \mathfrak A(\varphi)\le1,\qquad
 d_{Q_n}=(1-\varepsilon_n)\varepsilon_n^{-1/2},\qquad
 d_{(Q_n)_+}=0,\\
 \mathfrak D_\infty(\varphi)
 \ge\sup_n|d_{(Q_n)_+}-d_{Q_n}|
 =\sup_n(1-\varepsilon_n)2^n=\infty,\\
 \varphi\in\mathrm{bmo}_{\mu,2}\setminus\X_\infty(\mu),\qquad
 C_\varphi\notin\mathcal B(L^2(\mu)).
\end{gather*}
\end{remark}

\section{Comparison with established results}\label{sec-comparison}
We compare the operator, measure assumptions, and symbol condition
separately. The results attributed to the references below are
distinguished from consequences of Theorems~\ref{main}
and~\ref{boundedness} proved in this note.

\paragraph{The same downward shift and its paraproducts.}
L\'opez-S\'anchez, Martell and Parcet~\cite[Theorem 2.5]{LMP2014}
study the same operator $Hh_I=h_{I_-}-h_{I_+}$. With $q$ denoting an
$L^q$ exponent, their theorem gives
\[
 \begin{aligned}
 1<q<2:\qquad
 H\in\mathcal B(L^q(\mu))
  &\Longleftrightarrow
       \sup_{I\in\D}\frac{m(I)}{m(\widehat I)}<\infty,\\
 H^*\in\mathcal B(L^q(\mu))
  &\Longleftrightarrow
       \sup_{I\in\D}\frac{m(\widehat I)}{m(I)}<\infty.
 \end{aligned}
\]
The $L^2$ boundedness of $H$ is automatic. Their
Lemma~5.10 gives the Carleson embedding bound used here; testing
on $\mathbf1_Q/\sqrt{\mu(Q)}$ supplies its converse:
\[
 \mathfrak A(\varphi)^2
    =\sup_Q\frac1{\mu(Q)}\sum_{I\subseteq Q}|b_I|^2,\qquad
 \mathfrak A(\varphi)\le\norm{\pi_\varphi}
                         \le2\mathfrak A(\varphi).
\]
Thus the paraproduct boundedness ingredient is an established result.
Theorem~\ref{boundedness} concerns the additional symbol information
needed for the commutator on $L^2(\mu)$:
\[
 \norm{C_\varphi}\asymp
          \mathfrak A(\varphi)+\mathfrak D_\infty(\varphi).
\]

\paragraph{Martingale BMO and the diagonal condition.}
Treil~\cite[Proposition 5.1 and Theorem 5.3]{Treil2013} proves BMO
sufficiency for commutators with bounded martingale transforms and
a converse under mixing hypotheses, with the stated corrections
for finite-measure tail atoms. Such transforms preserve martingale
difference blocks. Here $H$ changes the dyadic level, so that theorem
does not directly identify its commutator symbols.

To compare the sufficient BMO condition with our exact condition, set
\[
 \begin{aligned}
 J(\varphi)&=\sup_{I\in\D}|\varphi_I-\varphi_{\widehat I}|,\\
 \mathfrak B(\varphi)&=
 \sup_{I\in\D}\left(\frac1{\mu(I)}
               \int_I|\varphi-\varphi_{\widehat I}|^2d\mu\right)^{1/2}.
 \end{aligned}
\]
The space $\{\mathfrak B<\infty\}$ is martingale
$\mathrm{BMO}_{\mu,2}$, whereas $\{\mathfrak A<\infty\}$ is
$\mathrm{bmo}_{\mu,2}$. Directly,
\begin{gather*}
 \frac1{\mu(I)}\int_I|\varphi-\varphi_{\widehat I}|^2d\mu
 =\frac1{\mu(I)}\int_I|\varphi-\varphi_I|^2d\mu
                     +|\varphi_I-\varphi_{\widehat I}|^2,\\
 \max\{\mathfrak A,J\}\le\mathfrak B
                   \le(\mathfrak A^2+J^2)^{1/2},\\
 |d_I-\varphi_I|
  \le\max\{|\varphi_{I_-}-\varphi_I|,
                   |\varphi_{I_+}-\varphi_I|\}\le J,\\
 |d_I-d_{\widehat I}|\le3J,\qquad
 \mathfrak D_\infty\le3\sqrt2J,\\
 \norm{\varphi}_{\X_\infty}
  \le\mathfrak A+3\sqrt2J\le(1+3\sqrt2)\mathfrak B,\\
 \mathrm{BMO}_{\mu,2}\subseteq\X_\infty(\mu)
                         \subseteq\mathrm{bmo}_{\mu,2}.
\end{gather*}
These inclusions are asserted without a strictness claim for the
first inclusion. The last example of the preceding section shows
that $\mathfrak A<\infty$ alone need not imply boundedness.

\paragraph{The sibling shift and the closest symbol characterization.}
Let $I^s$ denote the sibling of $I$ and define $\mathcal S$ by
\[
 \mathcal S h_{I_-}=-h_{I_+},\qquad
 \mathcal S h_{I_+}=h_{I_-},\qquad
 \mathcal S|_{\mathcal H_0}=0.
\]
Borges, Conde-Alonso, Pipher and Wagner~\cite[Theorem A]{BCAPW2025}
prove BMO upper bounds for this sibling shift and show that a
commutator lower bound for the full BMO norm fails in general.
D'Emilio, Lin, Wagner and Wick~\cite[Theorem D; Theorem 4.7]{DELWW2025}
give its exact symbol characterization. In the atomless,
sibling-balanced setting, with infinite measure on both dyadic
components, the $L^2$ specialization is
\begin{gather*}
 [\mu]_{\mathrm{sib}}
    :=\sup_{I\in\D}\frac{m(I)}{m(I^s)}<\infty,\qquad
 \beta_I=d_I-d_{I^s},\\
 [M_\varphi,\mathcal S]\in\mathcal B(L^2(\mu))
 \ \Longleftrightarrow\
 \mathfrak A(\varphi)<\infty,\quad\beta\in\ell^\infty.
\end{gather*}
Thus the use of $d_I=\langle\varphi,h_I^2\rangle$ already occurs
in that work. The respective differences follow the respective
shift actions:
\[
 \begin{aligned}
 \norm{[D_\varphi,H]}
    &=\sup_I\big(|d_{I_-}-d_I|^2+|d_{I_+}-d_I|^2\big)^{1/2},\\
 \norm{[D_\varphi,\mathcal S]}
    &=\sup_I|d_I-d_{I^s}|.
 \end{aligned}
\]
\begin{center}
\small
\begin{tabularx}{\linewidth}{@{}>{\raggedright\arraybackslash}p{0.24\linewidth}c
 >{\raggedright\arraybackslash}p{0.18\linewidth}
 >{\raggedright\arraybackslash}X@{}}
\toprule
Reference & Shift & Measure assumptions & Boundedness result on $L^2$\\
\midrule
Borges et al.\ \cite{BCAPW2025}
 & $\mathcal S$ & Sibling balanced
 & $\norm{[M_\varphi,\mathcal S]}\lesssim_{\mathrm{sib}}\mathfrak B(\varphi)$;
   full BMO lower bound fails.\\
D'Emilio et al.\ \cite{DELWW2025}
 & $\mathcal S$ & Atomless, sibling balanced
 & $\mathfrak A(\varphi)<\infty$ and
   $\sup_I|d_I-d_{I^s}|<\infty$.\\
Theorem~\ref{boundedness}
 & $H$ & $\mu\in\mathfrak M$
 & $\mathfrak A(\varphi)+\mathfrak D_\infty(\varphi)<\infty$.\\
\bottomrule
\end{tabularx}
\end{center}
The sibling-shift rows are read with $\mu(E_-)=\mu(E_+)=\infty$.
The subscript $\mathrm{sib}$ allows dependence on $[\mu]_{\mathrm{sib}}$.
The table compares distinct shifts. It does not identify the two
commutator symbol spaces for a fixed measure.

\paragraph{Schatten--Besov results and the scope of our criterion.}
For regular dyadic martingales with Lebesgue measure, the classical
Chao--Peng paraproduct theorem~\cite{CP1996}, recalled in
\cite[Theorem 1.1]{WZ2024}, gives a Besov characterization for every
$0<p<\infty$. The norm equivalence also follows from the scalar,
$d=2$ case of \cite[Theorem 1.2]{WZ2024}:
\begin{gather*}
 \norm{\varphi}_{B_p^{\D}}
    :=\left(\sum_{I\in\D}|I|^{-p/2}|b_I|^p\right)^{1/p},\\
 d\mu=dx,\quad0<p<\infty:
 \qquad \norm{\pi_\varphi}_{S_p}
    =\mathfrak C_p(\varphi)\asymp_p\norm{\varphi}_{B_p^{\D}}.
\end{gather*}
Wei and Zhang's theorem also treats operator-valued symbols on
regular $d$-adic filtrations. Our symbols are scalar, while the
measure is arbitrary. For the common scalar Lebesgue case,
\begin{gather*}
 \mathfrak D_p(\varphi)=\sqrt2\norm{\varphi}_{B_p^{\D}},\\
 1<p<\infty:\qquad
 \norm{C_\varphi}_{S_p}\asymp_p\norm{\varphi}_{B_p^{\D}}
 \qquad\text{by Theorem~\ref{main}}.
\end{gather*}
For arbitrary $\mu$, $\mathfrak C_p$ retains the interval correlations
through $G_F$, and Theorem~\ref{main} has the range $1<p<\infty$.
The atomic example in Section~\ref{sec-endpoint-example} prevents extending this separated
criterion to $p\le1$ for all measures; it does not conflict with the
regular-filtration paraproduct theorem. Theorem~\ref{all-p}, valid
for all $p>0$, is a finite-matrix reformulation and supplies no
coefficient-only Besov characterization in that endpoint range.

\begin{thebibliography}{9}
\bibitem{PST}
J. Parcet, M. de la Salle and E. Tablate,
\emph{The local geometry of idempotent Schur multipliers},
Forum Math. Pi \textbf{13} (2025), e14,
\href{https://doi.org/10.1017/fmp.2025.6}{doi:10.1017/fmp.2025.6},
\S 2.4.
\bibitem{LMP2014}
L. D. L\'opez-S\'anchez, J. M. Martell and J. Parcet,
\emph{Dyadic harmonic analysis beyond doubling measures},
Adv. Math. \textbf{267} (2014), 44--93,
\href{https://doi.org/10.1016/j.aim.2014.08.001}{doi:10.1016/j.aim.2014.08.001}.
\bibitem{Treil2013}
S. Treil,
\emph{Commutators, paraproducts and BMO in non-homogeneous martingale settings},
Rev. Mat. Iberoam. \textbf{29} (2013), no. 4, 1325--1372,
\href{https://doi.org/10.4171/RMI/759}{doi:10.4171/RMI/759}.
\bibitem{BCAPW2025}
T. Borges, J. M. Conde-Alonso, J. Pipher and N. A. Wagner,
\emph{Commutator estimates for Haar shifts with general measures},
J. Funct. Anal. \textbf{289} (2025), no. 5, Article 110945,
\href{https://doi.org/10.1016/j.jfa.2025.110945}{doi:10.1016/j.jfa.2025.110945}.
\bibitem{DELWW2025}
F. D'Emilio, Y. Lin, N. A. Wagner and B. D. Wick,
\emph{Optimal sparse bounds and commutator characterizations without doubling},
preprint (2025),
\href{https://arxiv.org/abs/2510.26505v1}{arXiv:2510.26505v1}.
\bibitem{CP1996}
J.-A. Chao and L. Peng,
\emph{Schatten classes and commutators on simple martingales},
Colloq. Math. \textbf{71} (1996), no. 1,
\href{https://doi.org/10.4064/cm-71-1-7-11}{doi:10.4064/cm-71-1-7-11}.
\bibitem{WZ2024}
Z. Wei and H. Zhang,
\emph{On the boundedness and Schatten class property of noncommutative
martingale paraproducts and operator-valued commutators},
preprint (2024),
\href{https://arxiv.org/abs/2407.14988v3}{arXiv:2407.14988v3}.
\end{thebibliography}
\end{document}
